\section{FORMAL RESTATEMENT}
\textbf{Erd\H{o}s \#29; reconstruction, 2026-09-23.}
Here $\mathbb N_0=\{0,1,\ldots\}$. Construct explicitly
$A\subseteq\mathbb N_0$ such that $A+A=\mathbb N_0$ and the ordered
representation function
\[
 r_A(n)=|\{(a,a')\in A^2:a+a'=n\}|
\]
satisfies $r_A(n)=o(n^\epsilon)$ for every $\epsilon>0$.
We give a membership rule computable in time polynomial in the number
of binary digits of its input and prove $r_A(n)=n^{o(1)}$.

\section{QUICK LITERATURE AND CONTEXT CHECK}
The problem is marked proved on its current page. Jain, Pham, Sawhney
and Zakharov gave an explicit construction in \emph{An explicit
economical additive basis} (2024), using a modular construction of
Ruzsa and mixed radix expansions. We reconstruct that method, including
the finite-field calculation behind the modular cover. The proof uses
the standard prime number theorem in the fixed progression $3$ modulo
$8$ only to control the sizes of the primes and the membership running
time. Neither the construction nor the conclusion is claimed as new.

\section{WORK}
\subsection{1. A modular basis with a uniform representation bound}
Let $p>12$ be a prime with $p\equiv3\pmod8$, so $2$ is a quadratic
nonresidue modulo $p$. For $c\in\{3,4,6\}$ and $0\leq x<p$ define
\[
 f_c(x)=x+2p\,[cx^2]_p,\qquad
 D_p=\{f_c(x):c\in\{3,4,6\},\ 0\leq x<p\},
\]
where $[z]_p$ is the representative in $\{0,\ldots,p-1\}$.
Let $D'_p=D_p+\{-p,0\}$ and let $E_p$ be its image modulo $p^2$,
represented by integers from $0$ to $p^2-1$.
We prove that for every residue $n$,
\[
 1\leq r_{E_p}(n\bmod p^2)\leq M,\qquad M=432.      \tag{1}
\]
In particular $0\in E_p$.

For coverage, write $0\leq n<p^2$ as $n=r+pq$, with $0\leq r,q<p$,
and put $t=q/2$ in $\mathbb F_p$. We seek $x+y=r$ and
$cx^2+dy^2=t$ in that field. Eliminating $y$ gives a quadratic in $x$
whose discriminant is
\[
 4\big((c+d)t-cd r^2\big).
\]
The choices $(c,d)=(3,6)$ and $(4,4)$ give respectively
$36(t-2r^2)$ and $32(t-2r^2)$. If $t-2r^2=0$ both quadratics have
a root. Otherwise the quotient of the discriminants is $8/9$, a
quadratic nonresidue, so precisely one discriminant is a square.
The leading coefficients $9$ and $8$ are nonzero modulo $p$.
Hence one of the two equations has a solution.

Choose representatives $x,y$ in $[0,p-1]$. Then $x+y=r+ep$ for some
$e\in\{0,1\}$, and
\[
 f_c(x)+f_d(y)\equiv n+ep\pmod{p^2}.
\]
Subtracting $ep$ from the first summand is allowed in $D'_p$, so every
residue is a sum of two elements of $E_p$.

For the upper bound, fix an integer $v$. In a representation
$v=f_c(x)+f_d(y)$, the value $s=x+y$ is uniquely specified by
$s\equiv v\pmod{2p}$ and $0\leq s\leq2p-2$, if it exists at all.
The congruence
\[
 cx^2+d(s-x)^2\equiv\frac{v-s}{2p}\pmod p
\]
has at most two roots for each ordered pair $(c,d)$: its leading
coefficient $c+d$ is between $6$ and $12$, and is nonzero modulo $p$.
Thus the number of ordered representations in $D_p$ is at most $18$.
There may be duplicate descriptions of an element by $(c,x)$; counting
all descriptions only increases this upper bound.

The two possible shifts for each summand imply $r_{D'_p}(v)\leq72$.
As $D'_p\subset[-p,2p^2)$, its pair sums lie in $[-2p,4p^2)$.
There are at most six integers in this interval with any fixed residue
modulo $p^2$. Choose one lift in $D'_p$ for each element of $E_p$;
this injects ordered residue representations into these integer
representations. The bound $6\cdot72=432$ proves (1).

\subsection{2. The explicit mixed radix set}
Let $p_i$ be the least prime congruent to $3$ modulo $8$ with
$p_i\geq i+13$, and put $b_i=p_i^2$. Repeated primes are permitted;
the bases are nondecreasing. The prime number theorem for this fixed
arithmetic progression gives $p_i=O(i+13)$, whereas $p_i\geq i+13$
holds by construction. Let $E_i=E_{p_i}$ and
$Q_i=\prod_{j<i}b_j$, with $Q_1=1$.

Every nonnegative integer has a canonical finite expansion
$a=\sum_{i=1}^{\ell}a_iQ_i$, with $0\leq a_i<b_i$ and nonzero top
digit unless $a=0$. Define $A$ by requiring
\[
 a_i\in E_i\quad(1\leq i<\ell),
 \qquad 0\leq a_\ell<b_\ell,                         \tag{2}
\]
and include $0$. Thus the top digit is unrestricted, while every lower
digit comes from its modular basis. Since $0\in E_i$, allowing leading
zeros gives the same set as the canonical rule.

\subsection{3. Every integer is represented, with carries checked}
For $n>0$, write its digits as $n_1,\ldots,n_k$ with $n_k\geq1$.
Start with carry $c_0=0$. For $i<k$, use (1) to choose
$a_i,a'_i\in E_i$ satisfying
\[
 a_i+a'_i+c_{i-1}\equiv n_i\pmod{b_i}.
\]
The next carry is the integer determined by
\[
 a_i+a'_i+c_{i-1}=n_i+c_i b_i.
\]
It belongs to $\{0,1\}$ because the left side is between $0$ and
$2b_i-1$. Set $a_k=n_k-c_{k-1}\geq0$ and $a'_k=0$.
This satisfies the top equality with no further carry. Both resulting
integers satisfy (2), even if a top digit vanishes, and their sum is
exactly $n$. For $n=0$, use $0+0$. Therefore $A+A=\mathbb N_0$.

\subsection{4. Bounding all representations, not only the constructed one}
For a fixed $n$ of digit length $k$, orient a representation so that
the smaller digit length of its two summands is $\ell\leq k$.
Treat zero as having length one. The orientation costs a factor at
most two. At each level $i<\ell$, both digits are restricted by (2)
and the incoming carry is already determined. There are at most $M$
ordered choices, by (1). At level $\ell$ choose the top digit of the
shorter summand in at most $b_\ell$ ways. Its lower digits and hence
that entire summand are then fixed; the other summand is forced to
be $n-a$. We have therefore bounded every representation by
\[
 r_A(n)\leq2\sum_{\ell=1}^k b_\ell M^{\ell-1}
 \leq2b_k M^k.                                      \tag{3}
\]
Since $b_i\geq(i+13)^2$, a $k$-digit positive integer satisfies
\[
 \log n\geq\sum_{i<k}2\log(i+13)\gg k\log k.
\]
Thus $k=O(\log n/\log\log n)$; also $b_k=O((k+13)^2)$.
Taking logarithms in (3) yields
\[
 r_A(n)\leq\exp\!\left(O\!\left(\frac{\log n}{\log\log n}\right)\right)
 =n^{o(1)}.
\]
For any prescribed $\epsilon>0$, the exponent is eventually below
$\epsilon/2$, so $r_A(n)/n^\epsilon\to0$, as required.

\subsection{5. Why the construction is explicit}
For an input of bit length $L$, at most $L$ bases need to be generated.
The required primes have size $O(L+13)$, so even trial division and a
sequential search in the progression find them in time polynomial in
$L$. Every set $E_i$ is given by the displayed three quadratic formulas
and two shifts and has a list of length $O(p_i)$. Repeated division
finds the mixed radix digits, after which the finitely many membership
tests in (2) take polynomial time in $L$. This establishes the usual
strong meaning of explicitness, rather than just an unbounded search
through infinite candidate sets. \hfill$\square$

\section{VERIFICATION AND LIMITS}
The batch script checks (1) exhaustively for $p=19,43,59$, checks the
two-discriminant coverage argument in each of these fields, and tests
the mixed radix membership and carry construction for all $n\leq5000$.
These finite tests support the calculation; the arguments above prove
the assertions for all permitted primes and all integers.
The uniform constant $432$ is deliberately loose. No optimization or
improvement of the published result is claimed. The fixed-progression
prime-size estimate is the standard external number-theoretic input;
the finite modular and mixed radix arguments are proved here.

\section{FINAL STATUS}
\textbf{SOLVED (known affirmative construction reconstructed).}
The write-up provides a complete construction and proof with the
explicitly stated prime-distribution theorem as an input. It does not
claim novelty or an independent audit of the formal proof reported on
the problem page.

\section{SOURCES}
\url{https://www.erdosproblems.com/29}.
V.~Jain, H.~T.~Pham, M.~Sawhney and D.~Zakharov,
\emph{An explicit economical additive basis}, especially Section~2
and its Ruzsa construction:
\url{https://arxiv.org/abs/2405.08650}.
Checked 2026-09-23. Checks are in
\texttt{verification/solved\_batch\_26\_checks.py}.

This measures completion of the reconstruction using the stated
classical prime-distribution input, not originality.

\section{COMPLETION ESTIMATE}
\textbf{COMPLETION: 100\%}
